\فصل{مفاهیم پایه و مبانی نظری}

هدف از این فصل، ارائه پایه‌های علمی، نظری و فنی مفاهیمی است که در طراحی و پیاده‌سازی سامانه‌های هوشمند مبتنی بر مدل‌های زبانی\پانویس{Language Models} و جستجوی وب نقش اساسی ایفا می‌کنند. پیش از ورود به معماری دقیق سیستم پیشنهادی در فصول آتی، درک عمیق مکانیزم‌های سنتی بازیابی اطلاعات\پانویس{Information Retrieval}، بررسی دقیق محدودیت‌های مدل‌های زبانی، معماری نوین تولید افزوده شده با بازیابی\پانویس{Retrieval-Augmented Generation} (\لر{RAG}) و چالش‌های پیچیده تعامل با محیط پویای وب الزامی است. این فصل به عنوان پیش‌نیازی مستحکم برای درک منطق، الگوریتم‌ها و نوآوری‌های مطرح شده در ادامه پایان‌نامه عمل می‌کند.

\قسمت{موتورهای جستجوی سنتی}

پیش از ظهور مدل‌های زبانی هوشمند و سامانه‌های تولیدگر\پانویس{Generative Systems}، دسترسی به اطلاعات در وب عمدتاً از طریق موتورهای جستجوی سنتی\پانویس{Traditional Search Engines} انجام می‌پذیرفت. معماری موتورهای جستجوی وب معمولاً شامل خزنده‌های توزیع‌شده، نمایه‌ساز، ساخت نمایه معکوس و جزء جستجوگر است \cite{brin1998anatomy}. کارکرد این موتورها بر چهار پایه اصلی استوار است:

\شروع{شمارش}

\فقره \مهم{خزش\پانویس{Crawling} (Crawling):} ربات‌های نرم‌افزاری\پانویس{Software Bots} صفحات وب را به صورت گراف‌پیمایی\پانویس{Graph Traversal} پیمایش کرده و محتوای آن‌ها را برای پردازش‌های بعدی جمع‌آوری می‌کنند.

\فقره \مهم{نمایه‌سازی\پانویس{Indexing} (Indexing):} محتوای جمع‌آوری‌شده پس از پیش‌پردازش متنی\پانویس{Text Preprocessing} (نظیر حذف کلمات توقف\پانویس{Stop Words} و ریشه‌یابی\پانویس{Stemming}) در ساختارهای داده‌ای بهینه‌سازی‌شده مانند نمایه‌ی معکوس\پانویس{Inverted Index} (\لر{Inverted Index}) ذخیره می‌شود تا جستجو با پیچیدگی زمانی حداقل و سرعت بالا انجام پذیرد \cite{manning2008ir}.

\فقره \مهم{بازیابی و رتبه‌بندی\پانویس{Retrieval \& Ranking} (Retrieval \& Ranking):} هنگامی که کاربر پرس‌وجویی\پانویس{Query} ارسال می‌کند، سیستم با استفاده از الگوریتم‌های آماری نظیر $TF-IDF$ \cite{salton1988term} یا نسخه‌های پیشرفته‌تر مبتنی بر احتمالات نظیر $BM25$ \cite{robertson2009bm25}، مرتبط‌ترین اسناد را بازیابی می‌کند. سپس نتایج بر اساس فاکتورهایی چون ارتباط معنایی کلیدواژه‌ها و اعتبار مبتنی بر ساختار لینک‌ها\پانویس{Link-based Authority} (مانند الگوریتم $PageRank$ \cite{page1999pagerank}) مرتب می‌شوند.

\فقره \مهم{نمایش نتایج:} خروجی فرآیند به صورت لیستی از پیوندها (\لر{Links}) و خلاصه‌های کوتاه\پانویس{Snippets} (\لر{Snippets}) به کاربر ارائه می‌شود.

\پایان{شمارش}

اگرچه این موتورها در یافتن صفحات وب بسیار مقیاس‌پذیر\پانویس{Scalable} و کارآمد هستند، اما عمدتاً وظیفه‌ی یافتن و رتبه‌بندی اسناد را برعهده دارند و معمولاً پاسخ نهایی را به صورت مستقیم تولید نمی‌کنند. بنابراین برای پرسش‌های تحلیلی که نیازمند ترکیب اطلاعات از چندین منبع متمایز و استنتاج منطقی\پانویس{Logical Inference} هستند، به تنهایی کافی نیستند \cite{metzler2021rethink}. این محدودیت بنیادین، نقطه شروع گذار پژوهشگران به سمت سامانه‌های هوشمندتر مبتنی بر هوش مصنوعی بود.

\قسمت{مدل‌های زبانی بزرگ\پانویس{Large Language Models} \لر{(LLMs)} و پاسخ‌گویی محاوره‌ای}

مدل‌های زبانی بزرگ شبکه‌های عصبی عمیقی\پانویس{Deep Neural Networks} (غالباً مبتنی بر معماری ترانسفورمر\پانویس{Transformer Architecture} و مکانیزم توجه\پانویس{Attention Mechanism} \cite{vaswani2017attention}) هستند که روی حجم عظیمی از داده‌های متنی آموزش دیده‌اند. این مدل‌ها قادرند با پیش‌بینی توکن بعدی\پانویس{Next Token Prediction} با استفاده از احتمالات ریاضی، متن‌هایی بسیار روان، منسجم و انسان‌گونه تولید کنند. در مقیاس بزرگ، همین هدف آموزشی باعث ظهور توانایی‌هایی مانند یادگیری با چند نمونه\پانویس{Few-shot Learning} می‌شود \cite{brown2020language}. با این وجود، استفاده مستقیم از این مدل‌ها به عنوان موتور پاسخ‌گویی دارای چالش‌های جدی زیر است:

\شروع{فقرات}

\فقره \مهم{توهم\پانویس{Hallucination} (Hallucination):} مدل‌های زبانی گاهی پاسخ‌هایی تولید می‌کنند که از نظر قواعد گرامری کاملاً صحیح به نظر می‌رسند، اما از نظر حقایق بیرونی کاملاً ساختگی و غلط هستند. این توهم می‌تواند به شکل جعل واقعیت یا تقلید از باورهای نادرست انسانی بروز کند \cite{huang2025hallucination, lin2022truthfulqa}.

\فقره \مهم{دانش ایستا و قدیمی\پانویس{Static Knowledge} (Static Knowledge):} دانش جاسازی‌شده در وزن‌های مدل محدود به داده‌هایی است که تا زمان آموزش\پانویس{Knowledge Cut-off} دیده‌است. بنابراین مدل بدون جستجوی زنده، در مورد اخبار یا رویدادهای جدید گرایش به ارائه پاسخ‌های منسوخ دارد \cite{kasai2024realtimeqa}.

\فقره \مهم{نبود ارجاع‌پذیری\پانویس{Lack of Attribution} (Lack of Attribution):} مدل به صورت ذاتی نمی‌تواند نشان دهد که پاسخ تولید شده را از کدام منبع یاد گرفته است، که این امر اعتمادپذیری\پانویس{Reliability} سیستم را در کاربردهای حساس (پزشکی، حقوقی و آکادمیک) به شدت کاهش می‌دهد.

\فقره \مهم{محدودیت پنجره بافت\پانویس{Context Window} (Context Window):} مدل‌ها محدودیت حافظه در پردازش توکن‌ها دارند. افزون بر این، پدیده‌ای به نام «گم‌شدن در میانه»\پانویس{Lost in the Middle} نشان می‌دهد که حتی با افزایش طول پنجره بافت، مدل‌ها لزوماً از تمام اطلاعات به‌طور یکنواخت استفاده نمی‌کنند و تمرکز آن‌ها بیشتر بر ابتدا و انتهای متن است \cite{liu2024lost}.

\پایان{فقرات}

\قسمت{تولید افزوده شده با بازیابی\پانویس{Retrieval-Augmented Generation} \لر{(RAG)}}

برای غلبه بر چالش‌های ذاتی مدل‌های زبانی، معماری \مهم{تولید افزوده شده با بازیابی} معرفی شد \cite{lewis2020rag}. این رویکرد حافظه پارامتریک (وزن‌های مدل) را با دانش ماژولار و غیرپارامتریک (اسناد خارجی) ترکیب می‌کند \cite{guu2020realm}. معماری کلاسیک \لر{RAG} از دو بخش اصلی تشکیل شده است:

\شروع{شمارش}

\فقره \مهم{بازیاب\پانویس{Retriever} (Retriever):} ابتدا پرسش کاربر با استفاده از یک مدل جاسازی\پانویس{Embedding Model} به یک بردار ریاضی چگال\پانویس{Dense Vector} در فضای ابعاد بالا تبدیل می‌شود. سپس در یک پایگاه‌داده برداری\پانویس{Vector Database}، جستجوی همسایگی (مانند $Cosine \ Similarity$ یا روش‌های تقریبی $ANN$ \cite{johnson2019billion}) انجام شده و $Top-k$ سند مرتبط از نظر معنایی بازیابی می‌شود \cite{karpukhin2020dpr}.

\فقره \مهم{تولیدکننده\پانویس{Generator} (Generator):} اسناد بازیابی شده به عنوان «بافت خارجی\پانویس{External Context}» به همراه پرسش اصلی کاربر، در قالب یک پرامپت ساختاریافته به \لر{LLM} تزریق می‌شوند. مدل زبانی اکنون موظف است پاسخ را منحصراً بر اساس این اطلاعات ارائه‌شده تولید کند.

\پایان{شمارش}

این رویکرد می‌تواند پدیده توهم را به شکل قابل‌توجهی کاهش داده و امکان ارجاع‌پذیری را فراهم کند. با این حال، کیفیت پاسخ نهایی مستقیماً به کیفیت بازیابی و میزان تبعیت مدل از شواهد ارائه‌شده بستگی دارد.

\قسمت{تفاوت \لر{RAG} کلاسیک با \لر{RAG} مبتنی بر وب (چالش‌های اساسی)}

معماری \لر{RAG} در ابتدا روی مجموعه‌های متنی ایستا و ساختاریافته (مانند مقالات ویکی‌پدیا) آزمایش شد. اما پیاده‌سازی این الگو بر روی گستره‌ی نامتناهی «وب» چالش‌های فنی متفاوتی را به همراه دارد:

\شروع{فقرات}

\فقره \مهم{پویایی و مقیاس‌پذیری\پانویس{Dynamism and Scalability}:} محیط وب مجموعه‌ای عظیم و دائماً در حال تغییر است. بنابراین ساخت و نگهداری یک نمایه برداری کامل و محلی از کل وب برای سامانه‌های محلی عملی نیست.

\فقره \مهم{ناهمگنی، نویز و اسپم\پانویس{Heterogeneity and Spam}:} صفحات وب دارای ساختارهای متغیر، تبلیغات، قالب‌های تکراری و اطلاعات نامعتبر هستند. استخراج متن خالص و حذف نویزها\پانویس{Boilerplate Removal} \cite{kohlschutter2010boilerplate} در کنار مقابله با محتوای فریبنده و اسپم \cite{castillo2010adversarial} یک ضرورت پیچیده در این محیط است.

\فقره \مهم{نیاز به تفکیک پرسش‌ها\پانویس{Query Decomposition} (Query Decomposition):} در محیط وب، پاسخ به یک سؤال تحلیلی معمولاً در یک سند واحد یافت نمی‌شود \cite{yang2018hotpotqa}. سیستم هوشمند باید بتواند یک سؤال پیچیده را به چند زیرسؤال شکسته و جستجوهای متعددی\پانویس{Multi-hop Retrieval} را انجام دهد.

\پایان{فقرات}

\قسمت{سامانه‌های مبتنی بر جستجوی وب \لر{(Search-Augmented LLMs)} و ابزارهای \لر{API}}

به منظور مقابله با چالش‌های یاد شده، سامانه‌های پیشرفته امروزی مدل‌های زبانی را مستقیماً به ابزارهای خارجی نظیر جستجوی وب (\لر{API}) متصل می‌کنند \cite{nakano2021webgpt, schick2023toolformer}. مراحل پردازشی در این سامانه‌ها عموماً به شرح زیر است:

\شروع{شمارش}

\فقره \مهم{بازنویسی پرس‌وجو\پانویس{Query Reformulation} (Query Reformulation):} از آنجا که پرسش طبیعی کاربر لزوماً بهترین عبارت برای موتور جستجو نیست، مدل زبانی آن را به مجموعه‌ای از کلیدواژه‌های بهینه تبدیل می‌کند \cite{ma2023query}.

\فقره \مهم{انتخاب منابع\پانویس{Source Selection}:} نتایج خام دریافتی از \لر{API} ارزیابی شده و سایت‌های نامعتبر یا نتایج نامربوط فیلتر می‌شوند.

\فقره \مهم{استدلال چندمرحله‌ای و بازیابی فعال\پانویس{Active Retrieval}:} مدل در چرخه‌ای از فکر و عمل \cite{yao2023react} متون را بررسی می‌کند. در صورت نقص اطلاعات، در طول فرآیند تولید پاسخ، پرسش‌های جدیدی تولید کرده و فاز جستجو را تکرار می‌کند \cite{jiang2023active}.

\فقره \مهم{تولید ارجاعات\پانویس{Citation Generation}:} در گام نهایی، مدل زبانی در حین تولید پاسخ، با استفاده از شناسه‌های یکتا مشخص می‌کند که هر ادعا از کدام پیوند استخراج شده است.

\پایان{شمارش}

\قسمت{نقش تاریخچه گفتگو در پاسخ‌گویی محاوره‌ای}

یکی از ارکان اساسی سامانه‌های دستیار هوشمند مدرن، حفظ حالت محاوره‌ای\پانویس{Conversational State} است. در یک گفتگوی چندنوبتی، کاربران تمایل دارند از ضمایر اشاره استفاده کنند یا سؤالات خود را با فرض آگاهی سیستم از مکالمات قبلی مطرح نمایند. 

بدون مدیریت دقیق تاریخچه، سؤال فعلی کاربر به تنهایی فاقد معنای کامل بوده و موتور بازیابی را گمراه می‌کند. در این سامانه‌ها، بازنویسی پرسش وابسته به گفتگو\پانویس{Conversational Query Rewriting} انجام می‌شود تا پرسش فعلی با استفاده از بافتار مکالمات پیشین به یک پرسش مستقل و شفاف برای موتور جستجو تبدیل گردد \cite{wu2022conqrr}.

\قسمت{بازرتبه‌بندی\پانویس{Re-ranking} یا \لر{Re-ranking}}

موتورهای جستجو و \لر{API}های وب به دلیل نیاز به پاسخ‌گویی در مقیاس میلی‌ثانیه، معمولاً از مدل‌های رمزگذار دوگانه\پانویس{Bi-Encoders} \cite{reimers2019sentence} یا تطابق دقیق کلمات استفاده می‌کنند. در این روش‌ها، پرسش و سند به صورت مستقل پردازش می‌شوند. خروجی اولیه این موتورها ممکن است صدها سند مرتبط را بازگرداند که از نظر معنایی دقیقاً پاسخگوی سؤال پیچیده کاربر نیستند.

برای حل این مشکل از معماری دو مرحله‌ای استفاده می‌شود. پس از بازیابی اولیه و سریع اسناد، مفهوم \مهم{رتبه‌بندی مجدد} با استفاده از مدل‌های رمزگذار متقاطع\پانویس{Cross-Encoders} مطرح می‌گردد \cite{nogueira2019passage}. این مدل‌های یادگیری عمیق، پرسش کاربر و متن سند را به صورت همزمان دریافت کرده و ارتباط را در سطح کلمه-به-کلمه محاسبه می‌کنند. در سامانه‌های عملی، ترکیب بازیابی سریع اولیه با بازرتبه‌بندی دقیق، توازنی ایده‌آل میان فراخوانی، دقت و هزینه محاسباتی ایجاد می‌کند \cite{thakur2021beir} تا تنها بهترین اسناد به \لر{LLM} ارسال گردند.

\قسمت{معیارهای ارزیابی سامانه‌ها\پانویس{System Evaluation Metrics}}

برای سنجش علمی و کمی عملکرد سامانه‌های \لر{RAG} مبتنی بر جستجو، ارزیابی باید به صورت تفکیک‌شده در دو سطح «موتور بازیابی» و «موتور تولید پاسخ» انجام شود \cite{es2023ragas}. مهم‌ترین معیارهای استاندارد ارزیابی عبارتند از:

\شروع{فقرات}

\فقره \مهم{دقت\پانویس{Precision} و فراخوانی\پانویس{Recall}:} معیارهای کلاسیک ارزیابی بازیابی. دقت نشان می‌دهد چه کسری از اسناد بازیابی شده مرتبط بوده‌اند و فراخوانی نشان‌دهنده کسر اسناد مرتبط بازیابی شده نسبت به کل اسناد مرتبط موجود است.

\فقره \مهم{نرخ برخورد\پانویس{Hit Rate} و میانگین رتبه متقابل\پانویس{Mean Reciprocal Rank} (MRR):} بررسی اینکه آیا در نتایج اولیه حداقل یک سند حاوی پاسخ طلایی وجود داشته است و در چه رتبه‌ای قرار دارد.

\فقره \مهم{وفاداری\پانویس{Faithfulness}:} بررسی اینکه آیا محتوای تولید شده، مستقیماً از اسناد ارائه‌شده قابل استنتاج است یا مدل دچار توهم شده است.

\فقره \مهم{صحت ارجاع\پانویس{Citation Accuracy}:} معیاری دقیق‌تر برای اطمینان از اینکه ارجاع و پیوند ثبت شده در کنار هر خط از پاسخ، دقیقاً همان ادعای متنی را پشتیبانی می‌کند \cite{gao2023citations}.

\فقره \مهم{سودمندی\پانویس{Answer Relevance}:} بررسی میزان تطابق پاسخ نهایی با نیت اصلی کاربر و رفع نیاز اطلاعاتی وی.

\پایان{فقرات}

این معیارهای کمی و کیفی پایه‌های نظری لازم برای طراحی آزمایش‌ها، پیکربندی بهینه هایپرپارامترها\پانویس{Hyperparameters} و اثبات کارایی سامانه‌ی پیشنهادی در فصول آتی را فراهم می‌آورند.